\section{第42届复赛第6题：带盘簧的双球转速仪}
\subsection{题干}
转速仪的竖直转轴由轴承 $A,B$ 约束；质量为 $m_1,m_2$ 的小球固定在长为 $2l$ 的细杆两端，细杆中心绕 $O$ 点水平轴转动，并受回复力矩 $k\theta$。题目要求讨论匀速转动条件、角动量与轴承受力，以及受冲击后的非线性振动和小振动近似。
\begin{figure}[htbp]\centering\begin{tikzpicture}[>=Stealth,scale=.82]
  \draw[very thick] (0,-3.1)--(0,3.1);
  \draw[thick] (-.35,-3.1)--(.35,-3.1); \draw[thick] (-.65,-3.35)--(.65,-3.35);
  \foreach \x in {-.6,-.3,0,.3,.6} {\draw[thick] (\x,-3.35)--(\x-.25,-3.6);}
  \draw[thick] (-.18,1.75) rectangle (.18,1.95); \node[right] at (.25,1.85) {$A$};
  \draw[thick] (-.18,-1.75) rectangle (.18,-1.95); \node[right] at (.25,-1.85) {$B$};
  \fill (0,0) circle (2pt) node[below left] {$O$};
  \draw[thick] (-2.2,1.45)--(2.2,-1.45);
  \fill (-2.2,1.45) circle (4pt) node[left] {$a$}; \fill (2.2,-1.45) circle (4pt) node[right] {$b$};
  \draw[->] (0,0)--(2.6,0) node[right] {$x$}; \draw[->] (0,0)--(0,3.5) node[above] {$z$};
  \draw[->] (-.5,3.1) arc (190:350:.5); \draw[->] (.5,3.1) arc (-10:170:.5);
  \draw (0,1.0) arc (90:146:1.0); \node at (-.45,.9) {$\theta$};
  \draw[<->] (.55,0)--(.55,1.85) node[midway,right] {$h$};
  \draw[<->] (.55,0)--(.55,-1.85) node[midway,right] {$h$};
  \node[below] at (0,-4.15) {原题图6a：带盘簧的双球转速仪};
\end{tikzpicture}
\caption{第42届复赛第6题原题图6a的TikZ重绘。}\end{figure}
\subsection{参考解答整理}
匀速转动的力矩平衡给出
\[\omega_0=\sqrt{\frac{k\theta_0-(m_1-m_2)gl\sin\theta_0}{(m_1+m_2+M/3)l^2\sin\theta_0\cos\theta_0}}.\]
相对 $O$ 点的总角动量为
\[\boldsymbol L=(m_1+m_2+M/3)l^2\sin\theta_0\,\omega_0(\cos\theta_0\boldsymbol e_x+\sin\theta_0\boldsymbol e_z).\]
轴承合力由质心加速度和角动量定理确定；$F_{Ax},F_{Bx}$ 的完整表达式见答案第18页，本整理不转抄未逐项复核的长式。

当 $m_1=m_2=m$ 时，竖直轴角动量守恒给出
\[\omega(\theta)=\frac{\sin^2\theta_0}{\sin^2\theta}\omega_0.\]
冲击后的初始角速度变化率满足
\[\Omega_0=\sqrt{\frac{k}{(2m+M/3)l^2}}\sqrt{\theta_0\tan\theta_0\left(\frac{\sin^2\theta_0}{\sin^2\theta_M}-1\right)+\theta_M^2-\theta_0^2}.\]
小幅振动频率为
\[\omega_\theta=\sqrt{\frac{k}{(2m+M/3)l^2}\left[1+\frac{\theta_0(2\cos^2\theta_0+1)}{\sin\theta_0\cos\theta_0}\right]}.\]
若 $\theta(t)=\theta_0+\delta_0\sin(\omega_\theta t)$，则一阶近似为
\[\omega(t)\approx\left[1-2\cot\theta_0\,\delta_0\sin(\omega_\theta t)\right]\omega_0.\]
\subsection{转写依据}
题干依据 \texttt{CPHO42-Q06\_transcription\_v1.md}；公式与分值依据答案第16至20页和 \texttt{CPHO42-Q06\_formula\_check\_v1.csv}。图6a由原题图重绘，不嵌入原始页面截图。
